% Standalone-preprint adaptation of cusp_compact_link_gaussian.tex.
% Canonical source SHA-256: 4469781B2386225453DBDD5C12B5B5D9213CD67480055F3F318143941C44FA52.
% The canonical radial-height symbol T was renamed R consistently; transpose
% notation was retained exactly. Bibliography keys were aligned with this preprint.

\section{Cubical holonomy and finite gauge systems}
\label{cusp-link-section}

This section records three constructions on the regular cusp fibres while
keeping two independent integers separate throughout:
\(M_{\mathrm{cov}}\) is the index used to enlarge two periods of a fibre,
whereas \(N_{\mathrm{mesh}}\) is the number of cubical subdivisions in each
coordinate direction.  In particular, no mesh refinement is identified with
a change of covering space.

Let
\begin{equation}
 P_R=
 \begin{pmatrix}
  6\operatorname{Re}\mu_R&\operatorname{Re}\tau_R&1&0\\
  6\operatorname{Im}\mu_R&\operatorname{Im}\tau_R&0&0\\
  \operatorname{Re}\beta_R&\operatorname{Re}\mu_R&0&1\\
  \operatorname{Im}\beta_R&\operatorname{Im}\mu_R&0&0
 \end{pmatrix},
 \qquad
 F_R=\mathbb R^4/P_R\mathbb Z^4.
 \label{cusp-link-period-matrix}
\end{equation}
Write
\begin{equation}
 L_R=\operatorname{Im}\tau_R,
 \qquad m_R=\operatorname{Im}\mu_R,
 \qquad q_R=-\operatorname{Im}\beta_R,
 \qquad D_R=L_Rq_R+6m_R^2>0.
 \label{cusp-link-radial-data}
\end{equation}
Expanding the determinant along the last two columns gives
\[
 \det P_R=L_R\operatorname{Im}\beta_R-6m_R^2=-D_R.
\]
Along the positive radial cusp the source data give
\begin{equation}
 L_R=R+O(1),\qquad q_R=R+O(1),\qquad m_R=O(1),
 \label{cusp-link-cusp-asymptotics}
\end{equation}
and all three real parts occurring in \eqref{cusp-link-period-matrix} are
bounded.  Equivalently, after ordering physical coordinates as
\((x_1,x_3;x_2,x_4)\), the first two columns have block form
\begin{equation}
\binom{\mathsf A_R^{\mathrm{blk}}}{\mathsf B_R^{\mathrm{blk}}},\qquad
\mathsf B_R^{\mathrm{blk}}=
\begin{pmatrix}6m_R&L_R\\-q_R&m_R\end{pmatrix}
=R\begin{pmatrix}0&1\\-1&0\end{pmatrix}
 +\mathsf C_R^{\mathrm{blk}},
\qquad \sup_R\lVert \mathsf C_R^{\mathrm{blk}}\rVert_{\mathrm{op}}<\infty.
 \label{cusp-link-block-form}
\end{equation}

For \(M_{\mathrm{cov}}\geq1\), put
\begin{equation}
 S_{M_{\mathrm{cov}}}
 =\operatorname{diag}(1,1,M_{\mathrm{cov}},M_{\mathrm{cov}}),
 \qquad
 P_{R,M_{\mathrm{cov}}}=P_RS_{M_{\mathrm{cov}}},
 \label{cusp-link-cover-matrix}
\end{equation}
and define the degree-\(M_{\mathrm{cov}}^2\) cover
\begin{equation}
 \widetilde F_{R,M_{\mathrm{cov}}}
 =\mathbb R^4/P_{R,M_{\mathrm{cov}}}\mathbb Z^4
 \longrightarrow F_R.
 \label{cusp-link-cover-torus}
\end{equation}
Indeed,
\(S_{M_{\mathrm{cov}}}\mathbb Z^4\subset\mathbb Z^4\) and
\[
 \mathbb Z^4/S_{M_{\mathrm{cov}}}\mathbb Z^4
 \cong(\mathbb Z/M_{\mathrm{cov}}\mathbb Z)^2,
\]
so the identity on \(\mathbb R^4\) induces the asserted cover and its deck
group has order \(M_{\mathrm{cov}}^2\).
Its Euclidean quotient metric, written on the coordinate torus
\(\mathbb R^4/\mathbb Z^4\), has Gram matrix and covolume
\begin{equation}
 G_{R,M_{\mathrm{cov}}}
 =P_{R,M_{\mathrm{cov}}}^{\mathsf T}P_{R,M_{\mathrm{cov}}},
 \qquad
 \mathcal D_{R,M_{\mathrm{cov}}}
 =\sqrt{\det G_{R,M_{\mathrm{cov}}}}
 =M_{\mathrm{cov}}^2D_R.
 \label{cusp-link-cover-metric}
\end{equation}

\subsection{The transported cubical subdivision and continuum holonomy}

Fix an integer \(N_{\mathrm{mesh}}\geq2\), and let
\(K_{N_{\mathrm{mesh}}}\) be the standard cubical subdivision of
\(\mathbb R^4/\mathbb Z^4\) obtained by dividing every coordinate circle
into \(N_{\mathrm{mesh}}\) equal intervals.  Its vertex set and positively
oriented coordinate edges are
\begin{equation}
 V_{N_{\mathrm{mesh}}}
 =(\mathbb Z/N_{\mathrm{mesh}}\mathbb Z)^4,
 \qquad
 E_{N_{\mathrm{mesh}}}
 =V_{N_{\mathrm{mesh}}}\times\{1,2,3,4\}.
 \label{cusp-link-cubical-sets}
\end{equation}

\begin{proposition}[Exact cubical realization and the torsion vertices]
\label{cusp-link-cubical-realization}
The linear map \(P_{R,M_{\mathrm{cov}}}\) induces a diffeomorphism
\begin{equation}
 \Phi_{R,M_{\mathrm{cov}}}:\mathbb R^4/\mathbb Z^4
 \longrightarrow\widetilde F_{R,M_{\mathrm{cov}}},
 \qquad [y]\longmapsto[P_{R,M_{\mathrm{cov}}}y].
 \label{cusp-link-coordinate-diffeomorphism}
\end{equation}
Transporting \(K_{N_{\mathrm{mesh}}}\) by this diffeomorphism gives an
honest cubical cell decomposition of the whole torus.  On vertices its
restriction is the bijection
\begin{equation}
 \begin{aligned}
 \iota_{R,M_{\mathrm{cov}},N_{\mathrm{mesh}}}:\quad
 &(\mathbb Z/N_{\mathrm{mesh}}\mathbb Z)^4
 \xrightarrow{\ \cong\ }
 \widetilde F_{R,M_{\mathrm{cov}}}[N_{\mathrm{mesh}}],\\
 &[n]\longmapsto
 \left[\frac1{N_{\mathrm{mesh}}}P_{R,M_{\mathrm{cov}}}n\right].
 \end{aligned}
 \label{cusp-link-torsion-bijection}
\end{equation}
Here the codomain is the finite \(N_{\mathrm{mesh}}\)-torsion subgroup,
not the entire torus.  The coordinate edge \((n,j)\) is the path
\begin{equation}
 \gamma_{n,j}(t)=
 \left[
  \frac1{N_{\mathrm{mesh}}}
  P_{R,M_{\mathrm{cov}}}(n+t\mathbf e_j)
 \right],
 \qquad 0\leq t\leq1.
 \label{cusp-link-coordinate-edge}
\end{equation}
More generally, for an ordered subset
\(I=\{i_1,\ldots,i_p\}\subset\{1,2,3,4\}\), the corresponding closed
\(p\)-cell is parametrized by
\begin{equation}
 c_{n,I}(t_1,\ldots,t_p)=
 \left[
  \frac1{N_{\mathrm{mesh}}}P_{R,M_{\mathrm{cov}}}
  \left(n+\sum_{a=1}^pt_a\mathbf e_{i_a}\right)
 \right],
 \qquad (t_1,\ldots,t_p)\in[0,1]^p.
 \label{cusp-link-general-cell}
\end{equation}
\end{proposition}

\begin{proof}
If \(y-y'\in\mathbb Z^4\), then
\(P_{R,M_{\mathrm{cov}}}(y-y')\) belongs to the quotient lattice, so
\eqref{cusp-link-coordinate-diffeomorphism} is well defined.  The matrix
\(P_{R,M_{\mathrm{cov}}}\) is invertible, and its inverse induces the
inverse map on the two quotients.  Hence it is a diffeomorphism and carries
every open cell of the standard subdivision diffeomorphically onto an open
cell of the target.

Every point displayed in \eqref{cusp-link-torsion-bijection} is killed by
\(N_{\mathrm{mesh}}\).  If it is zero, then
\[
 \frac{P_{R,M_{\mathrm{cov}}}n}{N_{\mathrm{mesh}}}
 \in P_{R,M_{\mathrm{cov}}}\mathbb Z^4.
\]
Invertibility gives
\(n\in N_{\mathrm{mesh}}\mathbb Z^4\), proving injectivity.  Conversely,
if \([x]\) is killed by \(N_{\mathrm{mesh}}\), choose a lift
\(x\in\mathbb R^4\).  There is \(n\in\mathbb Z^4\) such that
\(N_{\mathrm{mesh}}x=P_{R,M_{\mathrm{cov}}}n\), and therefore
\([x]=[P_{R,M_{\mathrm{cov}}}n/N_{\mathrm{mesh}}]\).  This proves
surjectivity.  Formula \eqref{cusp-link-coordinate-edge} is simply the
transport of the corresponding coordinate edge; at a wrap-around endpoint
the extra vector \(P_{R,M_{\mathrm{cov}}}\mathbf e_j\) is a period.
\end{proof}

Let \(G\) be a compact matrix Lie group and use the trivial principal bundle over
\(\widetilde F_{R,M_{\mathrm{cov}}}\).  A connection is represented by
\(A\in\Omega^1(\widetilde F_{R,M_{\mathrm{cov}}};\mathfrak g)\).  For an
oriented path \(\gamma\), define \(U_A(\gamma)\) by
\begin{equation}
 \frac{dU}{dt}=-A_{\gamma(t)}(\dot\gamma(t))U(t),
 \qquad U(0)=1,
 \qquad U_A(\gamma)=U(1).
 \label{cusp-link-parallel-transport-ode}
\end{equation}

\begin{proposition}[Typed continuum-to-link holonomy map]
\label{cusp-link-holonomy-map}
The formula
\begin{equation}
 \operatorname{Hol}_{R,M_{\mathrm{cov}},N_{\mathrm{mesh}}}(A)_{n,j}
 =U_A(\gamma_{n,j})
 \label{cusp-link-holonomy-definition}
\end{equation}
defines a map
\begin{equation}
 \operatorname{Hol}_{R,M_{\mathrm{cov}},N_{\mathrm{mesh}}}:
 \Omega^1(\widetilde F_{R,M_{\mathrm{cov}}};\mathfrak g)
 \longrightarrow G^{E_{N_{\mathrm{mesh}}}}.
 \label{cusp-link-holonomy-typed}
\end{equation}
For the convention
\begin{equation}
 A^g=g^{-1}Ag+g^{-1}dg,
 \qquad
 g\in C^\infty(\widetilde F_{R,M_{\mathrm{cov}}},G),
 \label{cusp-link-gauge-convention}
\end{equation}
it is equivariant for evaluation of \(g\) at the vertices:
\begin{equation}
 \operatorname{Hol}(A^g)_{n,j}
 =g(\gamma_{n,j}(1))^{-1}
  \operatorname{Hol}(A)_{n,j}
  g(\gamma_{n,j}(0)).
 \label{cusp-link-holonomy-equivariance}
\end{equation}
\end{proposition}

\begin{proof}
The smooth matrix ODE \eqref{cusp-link-parallel-transport-ode} has a unique
solution on every edge, which proves that \eqref{cusp-link-holonomy-typed}
is defined.  Let \(U\) solve that ODE for \(A\), set
\(g_t=g(\gamma_{n,j}(t))\), and put
\(V(t)=g_t^{-1}U(t)g_0\).  Direct differentiation gives
\[
 V'(t)=-\bigl(g_t^{-1}A(\dot\gamma_{n,j})g_t
                 +g_t^{-1}g_t'\bigr)V(t),
 \qquad V(0)=1.
\]
This is exactly the parallel-transport equation for \(A^g\).  Uniqueness
and evaluation at \(t=1\) prove \eqref{cusp-link-holonomy-equivariance}.
Notice that the endpoint inverse in this formula is forced by the joint
conventions \eqref{cusp-link-parallel-transport-ode} and
\eqref{cusp-link-gauge-convention}.
\end{proof}

\subsection{An explicit compact-link
\texorpdfstring{\(\mathrm{SU}(2)\)}{SU(2)} flux}

The ordered plaquette convention used below is most naturally written for
the inverse-transport link variable
\begin{equation}
 \mathcal U(A)_{n,j}
 =\operatorname{Hol}_{R,M_{\mathrm{cov}},N_{\mathrm{mesh}}}(A)_{n,j}^{-1}.
 \label{cusp-link-inverse-transport-variable}
\end{equation}
Equation \eqref{cusp-link-holonomy-equivariance} then becomes
\begin{equation}
 \mathcal U(A^g)_{n,j}
 =g(\gamma_{n,j}(0))^{-1}\mathcal U(A)_{n,j}g(\gamma_{n,j}(1)).
 \label{cusp-link-inverse-transport-gauge-law}
\end{equation}
Thus the product in \eqref{cusp-link-plaquette-convention} transforms by
conjugation at its base vertex.  The symbols \(U_j(n)\) in the remainder of
this subsection use this inverse-transport convention.

Use
\begin{equation}
 \mathfrak{su}(2)=\{X\in M_2(\mathbb C):X^*=-X,\ \operatorname{tr}X=0\},
 \qquad
 \langle X,Y\rangle_{\mathfrak{su}(2)}
 =-\frac12\operatorname{tr}(XY),
 \label{cusp-link-su2-inner-product}
\end{equation}
on real elements, and fix
\begin{equation}
 H_{\mathrm{link}}=\begin{pmatrix}i&0\\0&-i\end{pmatrix}.
 \qquad
 \lVert H_{\mathrm{link}}\rVert^2=1,
 \qquad \exp(2\pi H_{\mathrm{link}})=1.
 \label{cusp-link-flux-generator}
\end{equation}
Write \(N=N_{\mathrm{mesh}}\) temporarily inside the next displayed link
formulas, and represent \(n_1,n_2\) by integers in
\(\{0,\ldots,N-1\}\).

\begin{theorem}[A periodic compact-link flux and its exact action]
\label{cusp-link-flux-theorem}
On the oriented edges of \(K_N\), define
\begin{align}
 U_1(n)&=\exp\left(-\frac{2\pi n_2}{N^2}H_{\mathrm{link}}\right),
 \label{cusp-link-flux-u1}\\
 U_2(n)&=
 \begin{cases}
  1,&0\leq n_2\leq N-2,\\[1mm]
  \exp(2\pi n_1H_{\mathrm{link}}/N),&n_2=N-1,
 \end{cases}
 \label{cusp-link-flux-u2}\\
 U_3(n)&=U_4(n)=1.
 \label{cusp-link-flux-u34}
\end{align}
For
\begin{equation}
 U_{ij}(n)=U_i(n)U_j(n+\mathbf e_i)
            U_i(n+\mathbf e_j)^{-1}U_j(n)^{-1},
 \label{cusp-link-plaquette-convention}
\end{equation}
one has
\begin{equation}
 \begin{aligned}
 U_{12}(n)&=\exp(2\pi H_{\mathrm{link}}/N^2)
 &&\text{for every }n,\\
 U_{ij}(n)&=1
 &&\text{for }1\leq i<j\leq4,\quad(i,j)\neq(1,2).
 \end{aligned}
 \label{cusp-link-plaquette-values}
\end{equation}

Let \(\operatorname{Log}\) be the principal logarithm on the displayed
plaquette.  It is unambiguous because \(0<2\pi/N^2\leq\pi/2\).  Define
the logarithmic curvature two-form and its metric-weighted action by
\begin{align}
 \mathcal F_N^{\log}(n)
 &=N^2\operatorname{Log}(U_{12}(n))\,dy^1\wedge dy^2,
 \label{cusp-link-log-curvature}\\
 S^{\log}_{R,M_{\mathrm{cov}},N}(U)
 &=\frac{\mathcal D_{R,M_{\mathrm{cov}}}}{2g_{\mathrm{YM}}^2N^4}
   \sum_{n\in V_N}
   \lVert\mathcal F_N^{\log}(n)\rVert^2_{
       G_{R,M_{\mathrm{cov}}}^{-1}\otimes\mathfrak{su}(2)}.
 \label{cusp-link-log-action}
\end{align}
Then
\begin{equation}
 \mathcal F_N^{\log}(n)=2\pi H_{\mathrm{link}}\,dy^1\wedge dy^2,
 \qquad
 S^{\log}_{R,M_{\mathrm{cov}},N}(U)
 =\frac{2\pi^2M_{\mathrm{cov}}^2}{g_{\mathrm{YM}}^2D_R}.
 \label{cusp-link-log-action-value}
\end{equation}
In particular, the native-fibre value \(M_{\mathrm{cov}}=1\) is
\(2\pi^2/(g_{\mathrm{YM}}^2D_R)\), independently of the mesh.

There is no canonical scalar Wilson action for a general nonorthogonal
metric obtained merely by assigning independent character weights to all
coordinate plaquettes.  For this configuration, however, the exact
continuum-matched \((1,2)\)-sector Wilson functional is unambiguous if it is
defined explicitly by the following formula.  Its ordered plaquette
holonomy and real trace are the Wilson lattice-action construction
\cite{Wilson_1974}; the metric coefficient and normalization here are the
present definition, and their consequences are proved below rather than
imported from that source:
\begin{equation}
 S^{W,12}_{R,M_{\mathrm{cov}},N}(U)
 =\frac{\mathcal D_{R,M_{\mathrm{cov}}}\kappa^{12}_{R,M_{\mathrm{cov}}}}
        {2g_{\mathrm{YM}}^2}
   \sum_{n\in V_N}
   \left(2-\operatorname{Re}\operatorname{tr}U_{12}(n)\right),
 \label{cusp-link-wilson-sector-definition}
\end{equation}
where
\begin{equation}
 \kappa^{12}_{R,M_{\mathrm{cov}}}
 =(G_{R,M_{\mathrm{cov}}}^{-1})_{11}
  (G_{R,M_{\mathrm{cov}}}^{-1})_{22}
 -(G_{R,M_{\mathrm{cov}}}^{-1})_{12}^2
 =D_R^{-2}.
 \label{cusp-link-two-form-minor}
\end{equation}
Its exact value and continuum-mesh limit are
\begin{align}
 S^{W,12}_{R,M_{\mathrm{cov}},N}(U)
 &=\frac{M_{\mathrm{cov}}^2N^4}{2g_{\mathrm{YM}}^2D_R}
   \left(2-2\cos\frac{2\pi}{N^2}\right),
 \label{cusp-link-wilson-sector-value}\\
 &=\frac{M_{\mathrm{cov}}^2}{g_{\mathrm{YM}}^2D_R}
   \left(
    2\pi^2-\frac{2\pi^4}{3N^4}
    +\frac{4\pi^6}{45N^8}+O(N^{-12})
   \right),
 \label{cusp-link-wilson-sector-expansion}\\
 \lim_{N\to\infty}S^{W,12}_{R,M_{\mathrm{cov}},N}(U)
 &=\frac{2\pi^2M_{\mathrm{cov}}^2}{g_{\mathrm{YM}}^2D_R}.
 \label{cusp-link-wilson-sector-limit}
\end{align}
\end{theorem}

\begin{proof}
All links lie in the same diagonal circle subgroup, so they commute.  If
\(0\leq n_2\leq N-2\), the two \(U_2\)-terms in
\eqref{cusp-link-plaquette-convention} are trivial and
\[
 U_{12}(n)=
 \exp(-2\pi n_2H_{\mathrm{link}}/N^2)
 \exp(2\pi(n_2+1)H_{\mathrm{link}}/N^2)
 =\exp(2\pi H_{\mathrm{link}}/N^2).
\]
If \(n_2=N-1\), then \(U_1(n+\mathbf e_2)=1\).  Moreover,
\[
 U_2(n+\mathbf e_1)U_2(n)^{-1}=\exp(2\pi H_{\mathrm{link}}/N).
\]
For \(n_1=N-1\), the exponent computed with representatives differs from
\(2\pi H_{\mathrm{link}}/N\) by \(-2\pi H_{\mathrm{link}}\), whose exponential is one; hence the formula
also holds across that wrap boundary.  It follows that
\[
 U_{12}(n)=
 \exp\left(-\frac{2\pi(N-1)}{N^2}H_{\mathrm{link}}\right)
 \exp\left(\frac{2\pi}{N}H_{\mathrm{link}}\right)
 =\exp(2\pi H_{\mathrm{link}}/N^2).
\]
The functions \(U_1,U_2\) are independent of \(n_3,n_4\), and
\(U_3=U_4=1\), proving every remaining assertion in
\eqref{cusp-link-plaquette-values}.  Reversing the orientation gives
\(U_{21}(n)=U_{12}(n)^{-1}\); it is not an additional trivial plaquette.

The cover Gram matrix has lower-right principal block
\(M_{\mathrm{cov}}^2I_2\).  Jacobi's identity for complementary minors
therefore gives
\[
 \kappa^{12}_{R,M_{\mathrm{cov}}}
 =\frac{\det(M_{\mathrm{cov}}^2I_2)}
        {\det G_{R,M_{\mathrm{cov}}}}
 =\frac{M_{\mathrm{cov}}^4}{M_{\mathrm{cov}}^4D_R^2}
 =D_R^{-2},
\]
which proves \eqref{cusp-link-two-form-minor}.  By
\eqref{cusp-link-flux-generator} and the principal-log choice,
\(N^2\operatorname{Log}U_{12}=2\pi H_{\mathrm{link}}\).  Hence every one of the \(N^4\)
summands in \eqref{cusp-link-log-action} equals
\(4\pi^2D_R^{-2}\).  Substitution of
\(\mathcal D_{R,M_{\mathrm{cov}}}=M_{\mathrm{cov}}^2D_R\) proves
\eqref{cusp-link-log-action-value}.

Finally, \(H_{\mathrm{link}}\) has eigenvalues \(\pm i\), so
\[
 \operatorname{Re}\operatorname{tr}\exp(2\pi H_{\mathrm{link}}/N^2)
 =2\cos(2\pi/N^2).
\]
There are \(N^4\) equal plaquettes.  Substitution proves
\eqref{cusp-link-wilson-sector-value}.  Applying
\[
 1-\cos x=\frac{x^2}{2}-\frac{x^4}{24}
           +\frac{x^6}{720}+O(x^8)
\]
with \(x=2\pi/N^2\) proves
\eqref{cusp-link-wilson-sector-expansion} and
\eqref{cusp-link-wilson-sector-limit}.
\end{proof}

\subsection{Finite Gaussian quantization with independent cover and mesh}

Let
\begin{equation}
 \vartheta^j_{R,M_{\mathrm{cov}}}
 =P_{R,M_{\mathrm{cov}}}^{-\mathsf T}\mathbf e_j
 \in(\mathbb R^4)^*,
 \qquad 1\leq j\leq4.
 \label{cusp-link-physical-coframe}
\end{equation}
For \(0\leq p\leq4\), define the real finite cochain spaces
\begin{equation}
 C^p_{N_{\mathrm{mesh}}}(R,M_{\mathrm{cov}})
 =\operatorname{Map}\left(
 V_{N_{\mathrm{mesh}}},
 \bigwedge^p(\mathbb R^4)^*\otimes\mathfrak{su}(2)
 \right).
 \label{cusp-link-cochain-spaces}
\end{equation}
The exterior powers carry the metric induced by the physical coframe
\eqref{cusp-link-physical-coframe}; explicitly, for increasing multi-indices
\(I=(i_1,\ldots,i_p)\) and \(J=(j_1,\ldots,j_p)\),
\begin{equation}
 \left\langle
  \vartheta^{i_1}\wedge\cdots\wedge\vartheta^{i_p},
  \vartheta^{j_1}\wedge\cdots\wedge\vartheta^{j_p}
 \right\rangle
 =
 \det\left((G_{R,M_{\mathrm{cov}}}^{-1})_{i_a j_b}\right)_{a,b=1}^p.
 \label{cusp-link-exterior-metric}
\end{equation}
Together with \eqref{cusp-link-su2-inner-product}, this defines the
pointwise tensor-product metric.  The cochain inner product is
\begin{equation}
 \langle a,b\rangle_p
 =\frac{\mathcal D_{R,M_{\mathrm{cov}}}}{N_{\mathrm{mesh}}^4}
   \sum_{n\in V_{N_{\mathrm{mesh}}}}
   \langle a(n),b(n)\rangle.
 \label{cusp-link-cochain-inner-product}
\end{equation}
Let \(\mathsf T_j a(n)=a(n+\mathbf e_j)\), set
\begin{equation}
 D^+_{N_{\mathrm{mesh}},j}
 =N_{\mathrm{mesh}}(\mathsf T_j-1),
 \qquad
 D^-_{N_{\mathrm{mesh}},j}a(n)
 =N_{\mathrm{mesh}}(a(n)-a(n-\mathbf e_j)),
 \label{cusp-link-differences}
\end{equation}
and define
\begin{equation}
 d_{N_{\mathrm{mesh}}}a
 =\sum_{j=1}^4
  \varepsilon(\vartheta^j_{R,M_{\mathrm{cov}}})
  D^+_{N_{\mathrm{mesh}},j}a.
 \label{cusp-link-finite-exterior-derivative}
\end{equation}
The commuting differences and anticommuting exterior multiplications give
\(d_{N_{\mathrm{mesh}}}^2=0\).  Its adjoint is
\begin{equation}
 d_{N_{\mathrm{mesh}}}^*
 =-\sum_{j=1}^4
  \iota_{(\vartheta^j_{R,M_{\mathrm{cov}}})^\sharp}
  D^-_{N_{\mathrm{mesh}},j}.
 \label{cusp-link-finite-adjoint}
\end{equation}
Define the real reduced Coulomb space
\begin{align}
 \mathscr C_{R,M_{\mathrm{cov}},N_{\mathrm{mesh}}}
 &=
 \ker(d_{N_{\mathrm{mesh}}}^*:C^1\to C^0)
 \cap
 \left(\ker\Delta^{(1)}_{N_{\mathrm{mesh}}}\right)^\perp,
 \label{cusp-link-reduced-coulomb-space}\\
 \Delta^{(1)}_{N_{\mathrm{mesh}}}
 &=d_{N_{\mathrm{mesh}}}d_{N_{\mathrm{mesh}}}^*
 +d_{N_{\mathrm{mesh}}}^*d_{N_{\mathrm{mesh}}}.
 \label{cusp-link-finite-hodge-laplacian}
\end{align}
On this finite-dimensional space put
\begin{equation}
 \mathcal L_{R,M_{\mathrm{cov}},N_{\mathrm{mesh}}}
 =g_{\mathrm{YM}}^{-2}d_{N_{\mathrm{mesh}}}^*
                         d_{N_{\mathrm{mesh}}}.
 \label{cusp-link-finite-hessian}
\end{equation}

\begin{proposition}[Exact finite Hessian spectrum]
\label{cusp-link-finite-spectrum}
For
\begin{equation}
 \mathfrak d_{N_{\mathrm{mesh}}}(k)_j
 =N_{\mathrm{mesh}}
  \left(e^{2\pi i k_j/N_{\mathrm{mesh}}}-1\right),
 \qquad
 k\in(\mathbb Z/N_{\mathrm{mesh}}\mathbb Z)^4,
 \label{cusp-link-fourier-symbol}
\end{equation}
the positive spectrum of \(\mathcal L\) is obtained from
\begin{equation}
 \lambda_{R,M_{\mathrm{cov}},N_{\mathrm{mesh}}}(k)
 =\frac1{g_{\mathrm{YM}}^2}
  \left\lVert
   P_R^{-\mathsf T}S_{M_{\mathrm{cov}}}^{-1}
   \mathfrak d_{N_{\mathrm{mesh}}}(k)
  \right\rVert^2,
 \qquad k\neq0.
 \label{cusp-link-finite-mode-eigenvalue}
\end{equation}
Each complex Fourier fibre has transverse multiplicity
\(3\dim\mathfrak{su}(2)=9\).  In particular,
\begin{equation}
 \lambda_{\min}(\mathcal L)
 =\min_{k\neq0}\lambda_{R,M_{\mathrm{cov}},N_{\mathrm{mesh}}}(k)>0.
 \label{cusp-link-finite-minimum}
\end{equation}
For \(\mathbf e_1=(1,0,0,0)^{\mathsf T}\), the exact mode value is
\begin{equation}
 \lambda_{R,M_{\mathrm{cov}},N_{\mathrm{mesh}}}(\mathbf e_1)
 =\frac{4N_{\mathrm{mesh}}^2
              \sin^2(\pi/N_{\mathrm{mesh}})}{g_{\mathrm{YM}}^2}
  \frac{L_R^2+m_R^2}{D_R^2};
 \label{cusp-link-cover-soft-mode}
\end{equation}
it is independent of \(M_{\mathrm{cov}}\).
\end{proposition}

\begin{proof}
Translation is unitary for \eqref{cusp-link-cochain-inner-product}, so
\((D_j^+)^*=-D_j^-\).  Exterior multiplication by \(\vartheta^j\) has
adjoint contraction by \((\vartheta^j)^\sharp\), proving
\eqref{cusp-link-finite-adjoint}.  On the Fourier character
\(\chi_k(n)=e^{2\pi i k\cdot n/N_{\mathrm{mesh}}}\), the derivative symbol
is exterior multiplication by
\[
 \zeta(k)=P_{R,M_{\mathrm{cov}}}^{-\mathsf T}
           \mathfrak d_{N_{\mathrm{mesh}}}(k)
 =P_R^{-\mathsf T}S_{M_{\mathrm{cov}}}^{-1}
           \mathfrak d_{N_{\mathrm{mesh}}}(k).
\]
The Hermitian exterior-algebra identity
\[
 \iota_{\overline\zeta^\sharp}\varepsilon(\zeta)
 +\varepsilon(\zeta)\iota_{\overline\zeta^\sharp}
 =\lVert\zeta\rVert^2\operatorname{id}
\]
shows that on the Coulomb polarization
\(\ker\iota_{\overline\zeta^\sharp}\subset(\mathbb C^4)^*\), the operator
\(d^*d\) is multiplication by \(\lVert\zeta\rVert^2\).  For \(k\neq0\),
\(\mathfrak d_{N_{\mathrm{mesh}}}(k)\neq0\), so invertibility of
\(P_{R,M_{\mathrm{cov}}}\) gives
\(\zeta(k)\neq0\).  The contraction has rank one, and its kernel has
dimension three.  The zero Fourier fibre consists exactly of constant
harmonic cochains and was removed in
\eqref{cusp-link-reduced-coulomb-space}.  Fourier completeness now proves
\eqref{cusp-link-finite-mode-eigenvalue}--\eqref{cusp-link-finite-minimum}.

For \(k=\mathbf e_1\), the only nonzero component of
\(\mathfrak d_{N_{\mathrm{mesh}}}(k)\) has modulus
\[
 2N_{\mathrm{mesh}}\sin(\pi/N_{\mathrm{mesh}}),
\]
and
\(S_{M_{\mathrm{cov}}}^{-1}\mathbf e_1=\mathbf e_1\).  Solving
\(P_R^{\mathsf T}\xi=\mathbf e_1\), or directly inverting the lower
\(2\times2\) block in \eqref{cusp-link-period-matrix}, gives
\[
 \lVert P_R^{-\mathsf T}\mathbf e_1\rVert^2
 =\frac{L_R^2+m_R^2}{D_R^2}.
\]
Substitution proves \eqref{cusp-link-cover-soft-mode}.
\end{proof}

Let \(da\) denote Lebesgue measure on
\(\mathscr C_{R,M_{\mathrm{cov}},N_{\mathrm{mesh}}}\) induced by
\eqref{cusp-link-cochain-inner-product} and define
\begin{equation}
 \mathcal H^{\mathrm{Gauss}}_{R,M_{\mathrm{cov}},N_{\mathrm{mesh}}}
 =L^2(\mathscr C_{R,M_{\mathrm{cov}},N_{\mathrm{mesh}}},da).
 \label{cusp-link-gaussian-hilbert-space}
\end{equation}

\begin{theorem}[Finite Gaussian oscillator and a resolved cusp sequence]
\label{cusp-link-gaussian-theorem}
The Friedrichs realization on \eqref{cusp-link-gaussian-hilbert-space} of
\begin{equation}
 \widehat H^{\mathrm{Gauss}}_{R,M_{\mathrm{cov}},N_{\mathrm{mesh}}}
 =-\frac12\Delta_{\mathscr C}
  +\frac12\langle a,
     \mathcal L_{R,M_{\mathrm{cov}},N_{\mathrm{mesh}}}a\rangle
 \label{cusp-link-gaussian-hamiltonian}
\end{equation}
has a unique Gaussian ground state and spectral gap
\begin{equation}
 \operatorname{gap}(\widehat H^{\mathrm{Gauss}})
 =\sqrt{\lambda_{\min}(
   \mathcal L_{R,M_{\mathrm{cov}},N_{\mathrm{mesh}}})}.
 \label{cusp-link-gaussian-gap}
\end{equation}
Consequently,
\begin{equation}
 \operatorname{gap}(\widehat H^{\mathrm{Gauss}})
 \leq
 \frac{2N_{\mathrm{mesh}}\sin(\pi/N_{\mathrm{mesh}})}{g_{\mathrm{YM}}}
 \frac{\sqrt{L_R^2+m_R^2}}{D_R}.
 \label{cusp-link-gaussian-gap-bound}
\end{equation}

There is an explicit diagonal sequence for which the cover geometry tends locally to
\(\mathbb R^4\), the cubical mesh resolves that geometry, and the Gaussian
gap tends to zero.  Explicitly, for integers \(j\to\infty\), take
\begin{equation}
 R_j=j^2,\qquad M_{\mathrm{cov},j}=j,\qquad
 N_{\mathrm{mesh},j}=j^3.
 \label{cusp-link-diagonal-sequence}
\end{equation}
Then
\begin{equation}
 \begin{aligned}
 \operatorname{sys}(\widetilde F_{R_j,M_{\mathrm{cov},j}})
 &\longrightarrow\infty,\\
 \max_r
 \frac{\lVert P_{R_j,M_{\mathrm{cov},j}}\mathbf e_r\rVert}
      {N_{\mathrm{mesh},j}}
 &\longrightarrow0,\\
 \operatorname{gap}(\widehat H^{\mathrm{Gauss}})
 &\longrightarrow0,\\
 S^{\log}_{R_j,M_{\mathrm{cov},j},N_{\mathrm{mesh},j}}(U)
 &\longrightarrow0.
 \end{aligned}
 \label{cusp-link-diagonal-conclusions}
\end{equation}
Thus the pointed cover tori converge locally isometrically to \(\mathbb R^4\)
while their finite linearized Gaussian systems are simultaneously refined.
\end{theorem}

\begin{proof}
The real finite-dimensional operator \(\mathcal L\) is positive definite on
the reduced Coulomb space.  Choose an orthonormal eigenbasis with eigenvalues
\(0<\lambda_1\leq\cdots\leq\lambda_r\).  In the resulting coordinates,
\eqref{cusp-link-gaussian-hamiltonian} is the tensor sum
\[
 \sum_{s=1}^r
 \left(-\frac12\frac{\partial^2}{\partial x_s^2}
       +\frac12\lambda_sx_s^2\right).
\]
The normalized Hermite basis has eigenvalues
\(\sum_s(n_s+\tfrac12)\sqrt{\lambda_s}\), with
\(n_s\in\mathbb Z_{\geq0}\).  The all-zero multi-index is the unique ground
state, and the smallest possible positive excitation is
\(\sqrt{\lambda_1}\).  This proves \eqref{cusp-link-gaussian-gap}; applying
the minimum to the explicit \(\mathbf e_1\)-mode in
\eqref{cusp-link-cover-soft-mode} proves
\eqref{cusp-link-gaussian-gap-bound}.

For completeness, reorder physical coordinates as in
\eqref{cusp-link-block-form}.  Every vector of the cover lattice has the
form
\[
 \binom{\mathsf A_R^{\mathrm{blk}}n+M_{\mathrm{cov}}\ell}
       {\mathsf B_R^{\mathrm{blk}}n},
 \qquad n,\ell\in\mathbb Z^2.
\]
Let \(C_{\mathrm{blk}}=\sup_R
\lVert\mathsf C_R^{\mathrm{blk}}\rVert_{\mathrm{op}}\).  If \(n\neq0\), then
\(\lVert\mathsf B_R^{\mathrm{blk}}n\rVert
\geq(R-C_{\mathrm{blk}})\lVert n\rVert\); if \(n=0\) and the lattice
vector is nonzero, its norm is at least \(M_{\mathrm{cov}}\).  Hence
\begin{equation}
 \operatorname{sys}(\widetilde F_{R,M_{\mathrm{cov}}})
 \geq\min\{R-C_{\mathrm{blk}},M_{\mathrm{cov}}\}.
 \label{cusp-link-cover-systole-bound}
\end{equation}
This tends to infinity for the first two choices in
\eqref{cusp-link-diagonal-sequence}.  For any fixed \(r>0\), once the
systole exceeds \(4r\), the quotient map from \(\mathbb R^4\) is injective
and distance-preserving on the radius-\(r\) ball: for a nonzero period
\(\lambda\) and \(x,y\) in that ball,
\[
 \lVert x-y+\lambda\rVert
 \geq\lVert\lambda\rVert-\lVert x-y\rVert
 >4r-2r\geq\lVert x-y\rVert.
\]
This proves the pointed local-isometric \(\mathbb R^4\) limit.

The first two columns of \(P_R\) have norm \(O(R)\), while the last two
columns of \(P_{R,M_{\mathrm{cov}}}\) have norm exactly
\(M_{\mathrm{cov}}\).  Therefore the largest physical edge in the
\(N_{\mathrm{mesh}}\)-subdivision has length
\(O(\max\{R,M_{\mathrm{cov}}\}/N_{\mathrm{mesh}})\), which is \(O(j^{-1})\)
for \eqref{cusp-link-diagonal-sequence}.  Finally,
\eqref{cusp-link-cusp-asymptotics} gives
\(\sqrt{L_R^2+m_R^2}/D_R=R^{-1}(1+O(R^{-1}))\), while
\(2N\sin(\pi/N)\leq2\pi\).  The right-hand side of
\eqref{cusp-link-gaussian-gap-bound} is therefore \(O(j^{-2})\).
Moreover, \eqref{cusp-link-log-action-value} is \(O(j^{-2})\), since
\(M_{\mathrm{cov},j}^2/D_{R_j}=j^2/(j^4+O(j^2))\).
This proves every assertion in \eqref{cusp-link-diagonal-conclusions}.
\end{proof}

\subsection{The exponential map and the exact physical boundary}

Let \(C^1=C^1_{N_{\mathrm{mesh}}}(R,M_{\mathrm{cov}})\), regarded only as a
finite-dimensional Lie-algebra cochain space, and define componentwise
\begin{equation}
 \operatorname{Exp}_{N_{\mathrm{mesh}}}:C^1
 \longrightarrow\mathrm{SU}(2)^{E_{N_{\mathrm{mesh}}}},
 \qquad
 a\longmapsto\left(\exp(a_e/N_{\mathrm{mesh}})\right)_e.
 \label{cusp-link-exponential-map}
\end{equation}
Here, if
\(a(n)=\sum_{j=1}^4\vartheta^j_{R,M_{\mathrm{cov}}}\otimes a_j(n)\),
then \(a_e=a_j(n)\) for the oriented edge \(e=(n,j)\).

\begin{proposition}[Local chart, not a global embedding]
\label{cusp-link-exponential-local-only}
The map \eqref{cusp-link-exponential-map} is smooth and a local
diffeomorphism at zero, but it is not globally injective.  Under left
trivialization of every tangent space at the identity,
\begin{equation}
 D\operatorname{Exp}_{N_{\mathrm{mesh}}}(0)\,a
 =N_{\mathrm{mesh}}^{-1}a.
 \label{cusp-link-exponential-derivative}
\end{equation}
If \(S_W\) is any smooth compact-link action for which the identity
configuration is critical, then
\begin{equation}
 D^2(S_W\circ\operatorname{Exp}_{N_{\mathrm{mesh}}})_0[a,b]
 =N_{\mathrm{mesh}}^{-2}
  D^2(S_W)_{\mathbf1}[a,b]
 \label{cusp-link-exponential-hessian-scaling}
\end{equation}
under the same tangent identification.  Thus, for a prescribed cochain
quadratic form \(Q\),
\[
 D^2(S_W\circ\operatorname{Exp}_{N_{\mathrm{mesh}}})_0=Q
\]
holds exactly when
\[
 D^2(S_W)_{\mathbf1}=N_{\mathrm{mesh}}^2Q
\]
under the coefficient identification.  Both this scaling and the tensor
structure must be verified for the chosen compact-link action; neither
follows from the exponential map alone.
\end{proposition}

\begin{proof}
For the Lie-group exponential, \(D\exp_0\) is the identity.  Applying the
chain rule to \(a_e\mapsto a_e/N_{\mathrm{mesh}}\) gives
\eqref{cusp-link-exponential-derivative}, which is invertible, so the inverse
function theorem gives a local diffeomorphism.  It is not globally
injective: a cochain which is zero on every edge except one, where it equals
\(2\pi N_{\mathrm{mesh}}H_{\mathrm{link}}\), has the same image as the zero cochain by
\eqref{cusp-link-flux-generator}.  Finally, the second-derivative chain rule
contains a term involving \(DS_W(\mathbf1)\), which vanishes by criticality;
the two remaining first derivatives of \(\operatorname{Exp}\) each supply
the factor \(N_{\mathrm{mesh}}^{-1}\).  This proves
\eqref{cusp-link-exponential-hessian-scaling}.
\end{proof}

The Gaussian Hilbert space in \eqref{cusp-link-gaussian-hilbert-space} is
not the compact-link Yang--Mills Hilbert space.  At this finite lattice, the
latter kinematic and gauge-invariant spaces would be
\begin{equation}
 \mathcal H^{\mathrm{link}}_{N_{\mathrm{mesh}}}
 =L^2(\mathrm{SU}(2)^{E_{N_{\mathrm{mesh}}}},dU),
 \qquad
 \mathcal H^{\mathrm{inv}}_{N_{\mathrm{mesh}}}
 =\left(\mathcal H^{\mathrm{link}}_{N_{\mathrm{mesh}}}\right)^{
    \mathrm{SU}(2)^{V_{N_{\mathrm{mesh}}}}},
 \label{cusp-link-physical-hilbert-space}
\end{equation}
where \(dU\) is product normalized Haar measure and the superscript denotes
the nonlinear vertex-gauge action.  Here ``gauge-invariant'' refers only to
the link-configuration space of this four-dimensional cell complex.  Without
a chosen time direction or transfer matrix, it is not asserted to be a
canonical Hamiltonian physical Hilbert space.  L\"uscher's construction of
a selfadjoint, strictly positive transfer matrix
\cite{Luscher_1977} and Menotti--Pelissetto's reflection-positivity theorem
for the Wilson action \cite{MenottiPelissetto_1987} concern the standard
Wilson lattice setting after the requisite reflection/time structure has
been fixed; they do not supply those omitted choices for the nonorthogonal
sector functional \eqref{cusp-link-wilson-sector-definition}.  The
Osterwalder--Schrader reconstruction boundary is likewise the one stated in
the primary Euclidean axioms papers
\cite{OsterwalderSchrader_1973,OsterwalderSchrader_1975}, not a consequence of
a finite Gaussian oscillator calculation.  The differences are
structural:
\begin{enumerate}
\item \(\mathcal H^{\mathrm{Gauss}}\) is an \(L^2\)-space on a noncompact
linear Coulomb slice with Lebesgue measure; \(\mathcal H^{\mathrm{inv}}\)
is a gauge-invariant subspace of an \(L^2\)-space on compact link variables
with Haar measure.
\item The local exponential map is neither a global coordinate chart nor a
measure-preserving bijection.  No unitary map between the two displayed
Hilbert spaces has been constructed.
\item A link configuration \(U\) determines the Dirac distribution
\(\delta_U\), but
\begin{equation}
 \delta_U\in\mathcal D'(\mathrm{SU}(2)^{E_{N_{\mathrm{mesh}}}}),
 \qquad
 \delta_U\notin L^2(\mathrm{SU}(2)^{E_{N_{\mathrm{mesh}}}},dU).
 \label{cusp-link-dirac-not-l2}
\end{equation}
Indeed, product Haar measure is nonatomic, so a point mass is singular with
respect to it and has no square-integrable density.
\item No nonlinear compact-link Hamiltonian, Gauss-law spectral problem,
time direction, reflection-positive transfer matrix, renormalized-scale
limit, or Osterwalder--Schrader reconstruction is supplied here.
\end{enumerate}
Accordingly, \eqref{cusp-link-gaussian-gap} is an exact theorem about the
finite quadratic oscillator defined by the Coulomb-restricted operator
\(g_{\mathrm{YM}}^{-2}d^*d\).  It is not an eigenvalue statement about a
nonperturbative physical Yang--Mills
Hamiltonian, and the small classical action in
\eqref{cusp-link-log-action-value} is not by itself a physical-state Rayleigh
quotient.
